\chapter{Architecture 3 tiers - My Smart Budget}

\section{Architecture 3 tiers de My Smart Budget}

Dans cette section, je présente l'architecture 3 tiers de My Smart Budget, une application de gestion de budget personnel permettant aux utilisateurs de suivre leurs dépenses, gérer leurs budgets et analyser leurs habitudes financières.

\textbf{Architecture de My Smart Budget :} L'application suit une architecture 3 tiers stricte séparant la présentation (React), la logique métier (Node.js/Express) et les données (PostgreSQL/MongoDB).

\subsection{Couche Présentation - Interface de Gestion Budgétaire}

La couche présentation de My Smart Budget offre une interface intuitive pour la gestion des finances personnelles avec des tableaux de bord interactifs et des visualisations de données.

\textbf{Stack Frontend My Smart Budget :}

\begin{exemple}
\textbf{Technologies de présentation :}
\begin{itemize}
    \item \textbf{Framework :} React 18 avec TypeScript pour la sécurité des types
    \item \textbf{État :} Redux Toolkit pour gérer l'état des budgets et transactions
    \item \textbf{Routing :} React Router v6 pour la navigation entre modules
    \item \textbf{UI :} Material-UI + Chart.js pour les graphiques financiers
    \item \textbf{HTTP :} Axios avec intercepteurs pour l'authentification JWT
    \item \textbf{Validation :} Yup pour valider les montants et formulaires
\end{itemize}

\textbf{Structure des composants My Smart Budget :}
\begin{verbatim}
src/
+-- components/
|   +-- common/
|   |   +-- AmountInput.tsx        # Saisie de montants
|   |   +-- DateRangePicker.tsx    # Sélection de périodes
|   |   +-- CurrencyDisplay.tsx    # Affichage formaté
|   +-- dashboard/
|   |   +-- BudgetOverview.tsx     # Vue d'ensemble des budgets
|   |   +-- ExpenseChart.tsx       # Graphiques de dépenses
|   |   +-- AlertsPanel.tsx        # Alertes de dépassement
|   +-- transactions/
|   |   +-- TransactionList.tsx    # Liste des transactions
|   |   +-- TransactionForm.tsx    # Ajout/édition
|   |   +-- CategorySelector.tsx   # Sélection catégories
|   +-- budgets/
|   |   +-- BudgetCard.tsx         # Carte de budget
|   |   +-- BudgetProgress.tsx     # Barre de progression
|   |   +-- BudgetCreator.tsx      # Création de budget
|   +-- reports/
|       +-- MonthlyReport.tsx      # Rapport mensuel
|       +-- CategoryAnalysis.tsx   # Analyse par catégorie
|       +-- TrendsChart.tsx        # Tendances de dépenses
+-- hooks/
|   +-- useTransactions.ts         # Hook pour transactions
|   +-- useBudgetAlerts.ts        # Hook pour alertes
|   +-- useCurrency.ts            # Hook pour conversion
+-- services/
|   +-- budgetService.ts          # API des budgets
|   +-- transactionService.ts     # API des transactions
|   +-- analyticsService.ts       # API des analyses
\end{verbatim}
\end{exemple}

\subsection{Couche Logique Métier - Gestion Financière}

La couche métier de My Smart Budget implémente toute la logique de calcul budgétaire, de catégorisation automatique et d'analyse des dépenses.

\textbf{Architecture Backend My Smart Budget :}

\subsubsection{Controllers - Gestion des Requêtes Financières}

\textbf{Contrôleurs de My Smart Budget :}

\begin{exemple}
\textbf{TransactionController - Gestion des transactions :}
\begin{lstlisting}[language=JavaScript]
// TransactionController.js
class TransactionController {
  async createTransaction(req, res) {
    try {
      const { amount, category, description, date, type } = req.body;
      
      // Validation du montant et du type
      if (amount <= 0) {
        return res.status(400).json({ 
          error: 'Le montant doit être positif' 
        });
      }
      
      // Création de la transaction
      const transaction = await this.transactionService.create({
        userId: req.user.id,
        amount,
        category,
        description,
        date: date || new Date(),
        type // 'income' ou 'expense'
      });
      
      // Vérification du budget associé
      await this.budgetService.checkBudgetAlert(
        req.user.id, 
        category, 
        new Date()
      );
      
      res.status(201).json(transaction);
    } catch (error) {
      res.status(400).json({ error: error.message });
    }
  }
  
  async getMonthlyTransactions(req, res) {
    const { month, year } = req.params;
    const transactions = await this.transactionService
      .getByMonth(req.user.id, month, year);
    
    const summary = await this.analyticsService
      .calculateMonthlySummary(transactions);
    
    res.json({ transactions, summary });
  }
}
\end{lstlisting}
\end{exemple}

\subsubsection{Services - Logique Métier Financière}

\textbf{Services de My Smart Budget :}

\begin{exemple}
\textbf{BudgetService - Gestion des budgets :}
\begin{lstlisting}[language=JavaScript]
// BudgetService.js
class BudgetService {
  async createBudget(userId, budgetData) {
    // Validation des règles métier
    if (budgetData.amount <= 0) {
      throw new Error('Le montant du budget doit être positif');
    }
    
    // Vérification qu'un budget n'existe pas déjà
    const existingBudget = await this.budgetRepository
      .findByUserAndCategory(userId, budgetData.category, budgetData.period);
    
    if (existingBudget) {
      throw new Error('Un budget existe déjà pour cette catégorie');
    }
    
    // Création du budget avec calculs initiaux
    const budget = await this.budgetRepository.create({
      ...budgetData,
      userId,
      spent: 0,
      remaining: budgetData.amount,
      alertThreshold: budgetData.alertThreshold || 80 // 80% par défaut
    });
    
    // Configuration des alertes automatiques
    await this.notificationService.setupBudgetAlerts(budget);
    
    return budget;
  }
  
  async checkBudgetAlert(userId, category, date) {
    const budget = await this.budgetRepository
      .findActiveByCategory(userId, category, date);
    
    if (!budget) return;
    
    // Calcul du montant dépensé
    const spent = await this.transactionRepository
      .sumByCategory(userId, category, budget.startDate, budget.endDate);
    
    const percentageUsed = (spent / budget.amount) * 100;
    
    // Envoi d'alerte si dépassement du seuil
    if (percentageUsed >= budget.alertThreshold) {
      await this.notificationService.sendBudgetAlert({
        userId,
        budget,
        percentageUsed,
        remaining: budget.amount - spent
      });
    }
    
    // Mise à jour du budget
    await this.budgetRepository.updateSpent(budget.id, spent);
  }
}
\end{lstlisting}
\end{exemple}

\begin{exemple}
\textbf{AnalyticsService - Analyses financières :}
\begin{lstlisting}[language=JavaScript]
// AnalyticsService.js
class AnalyticsService {
  async generateMonthlyReport(userId, month, year) {
    // Récupération des données
    const transactions = await this.transactionRepository
      .getByMonth(userId, month, year);
    
    const budgets = await this.budgetRepository
      .getByMonth(userId, month, year);
    
    // Calculs des métriques
    const metrics = {
      totalIncome: this.calculateTotalByType(transactions, 'income'),
      totalExpenses: this.calculateTotalByType(transactions, 'expense'),
      savingsRate: 0,
      categoryBreakdown: this.groupByCategory(transactions),
      budgetAdherence: this.calculateBudgetAdherence(budgets, transactions),
      trends: await this.calculateTrends(userId, month, year)
    };
    
    metrics.savingsRate = ((metrics.totalIncome - metrics.totalExpenses) 
      / metrics.totalIncome) * 100;
    
    // Génération des insights
    const insights = await this.generateInsights(metrics, transactions);
    
    return {
      period: { month, year },
      metrics,
      insights,
      recommendations: this.generateRecommendations(metrics)
    };
  }
  
  generateRecommendations(metrics) {
    const recommendations = [];
    
    if (metrics.savingsRate < 20) {
      recommendations.push({
        type: 'savings',
        priority: 'high',
        message: 'Votre taux d\'épargne est inférieur à 20%',
        suggestion: 'Essayez de réduire vos dépenses non essentielles'
      });
    }
    
    // Analyse des catégories dépassant le budget
    metrics.categoryBreakdown.forEach(category => {
      if (category.percentageOfBudget > 100) {
        recommendations.push({
          type: 'overspending',
          category: category.name,
          priority: 'medium',
          message: `Dépassement de budget dans ${category.name}`,
          suggestion: `Réduisez de ${category.overspent}€ ce mois-ci`
        });
      }
    });
    
    return recommendations;
  }
}
\end{lstlisting}
\end{exemple}

\subsubsection{Repositories - Accès aux Données Financières}

\textbf{Repositories de My Smart Budget :}

\begin{exemple}
\textbf{TransactionRepository - Accès aux transactions :}
\begin{lstlisting}[language=JavaScript]
// TransactionRepository.js
class TransactionRepository {
  async create(transactionData) {
    return await prisma.transaction.create({
      data: {
        amount: transactionData.amount,
        type: transactionData.type,
        category: transactionData.category,
        description: transactionData.description,
        date: transactionData.date,
        userId: transactionData.userId,
        isRecurring: transactionData.isRecurring || false
      },
      include: {
        category: true,
        tags: true
      }
    });
  }
  
  async getByMonth(userId, month, year) {
    const startDate = new Date(year, month - 1, 1);
    const endDate = new Date(year, month, 0);
    
    return await prisma.transaction.findMany({
      where: {
        userId,
        date: {
          gte: startDate,
          lte: endDate
        }
      },
      orderBy: { date: 'desc' },
      include: {
        category: true,
        tags: true
      }
    });
  }
  
  async sumByCategory(userId, categoryId, startDate, endDate) {
    const result = await prisma.transaction.aggregate({
      where: {
        userId,
        categoryId,
        type: 'expense',
        date: {
          gte: startDate,
          lte: endDate
        }
      },
      _sum: {
        amount: true
      }
    });
    
    return result._sum.amount || 0;
  }
}
\end{lstlisting}
\end{exemple}

\subsection{Couche Données - Stockage des Données Financières}

My Smart Budget utilise une architecture de données hybride optimisée pour les données financières.

\textbf{Architecture de données My Smart Budget :}

\begin{exemple}
\textbf{Structure des données :}
\begin{itemize}
    \item \textbf{PostgreSQL :} Données transactionnelles (transactions, budgets, comptes)
    \item \textbf{MongoDB :} Logs d'audit, rapports générés, historique d'analyses
    \item \textbf{Redis :} Cache des calculs, sessions, taux de change
    \item \textbf{ORM :} Prisma pour PostgreSQL avec migrations versionnées
\end{itemize}

\textbf{Modèle de données PostgreSQL :}
\begin{lstlisting}[language=SQL]
-- Table des transactions
CREATE TABLE transactions (
  id UUID PRIMARY KEY DEFAULT gen_random_uuid(),
  user_id UUID NOT NULL REFERENCES users(id),
  amount DECIMAL(10,2) NOT NULL,
  type VARCHAR(10) CHECK (type IN ('income', 'expense')),
  category_id UUID REFERENCES categories(id),
  description TEXT,
  date TIMESTAMP NOT NULL,
  is_recurring BOOLEAN DEFAULT FALSE,
  recurring_pattern JSONB,
  tags TEXT[],
  created_at TIMESTAMP DEFAULT CURRENT_TIMESTAMP,
  updated_at TIMESTAMP DEFAULT CURRENT_TIMESTAMP
);

-- Table des budgets
CREATE TABLE budgets (
  id UUID PRIMARY KEY DEFAULT gen_random_uuid(),
  user_id UUID NOT NULL REFERENCES users(id),
  category_id UUID REFERENCES categories(id),
  amount DECIMAL(10,2) NOT NULL,
  period VARCHAR(20) CHECK (period IN ('monthly', 'weekly', 'yearly')),
  start_date DATE NOT NULL,
  end_date DATE NOT NULL,
  alert_threshold INTEGER DEFAULT 80,
  is_active BOOLEAN DEFAULT TRUE,
  created_at TIMESTAMP DEFAULT CURRENT_TIMESTAMP
);

-- Index pour optimiser les requêtes
CREATE INDEX idx_transactions_user_date ON transactions(user_id, date DESC);
CREATE INDEX idx_transactions_category ON transactions(category_id, date);
CREATE INDEX idx_budgets_user_active ON budgets(user_id, is_active);
\end{lstlisting}
\end{exemple}

\begin{exemple}
\textbf{MongoDB pour les analytics :}
\begin{lstlisting}[language=JavaScript]
// Collection des rapports mensuels
const monthlyReportSchema = {
  userId: ObjectId,
  period: {
    month: Number,
    year: Number
  },
  metrics: {
    totalIncome: Number,
    totalExpenses: Number,
    savingsRate: Number,
    categoryBreakdown: [{
      categoryId: String,
      categoryName: String,
      amount: Number,
      percentage: Number,
      trend: String // 'up', 'down', 'stable'
    }],
    dailySpending: [{
      date: Date,
      amount: Number
    }]
  },
  insights: [{
    type: String,
    severity: String,
    message: String,
    data: Object
  }],
  generatedAt: Date,
  version: String
};

// Pipeline d'agrégation pour les tendances
const spendingTrendsPipeline = [
  {
    $match: {
      userId: userId,
      'period.year': year
    }
  },
  {
    $group: {
      _id: '$period.month',
      totalExpenses: { $first: '$metrics.totalExpenses' },
      categories: { $first: '$metrics.categoryBreakdown' }
    }
  },
  {
    $sort: { _id: 1 }
  },
  {
    $project: {
      month: '$_id',
      totalExpenses: 1,
      topCategory: { $arrayElemAt: ['$categories', 0] },
      monthOverMonth: {
        $subtract: ['$totalExpenses', { $lag: '$totalExpenses' }]
      }
    }
  }
];
\end{lstlisting}
\end{exemple}

\subsection{Communication entre les tiers de My Smart Budget}

\textbf{Architecture de communication :}

\begin{exemple}
\textbf{Flux de données My Smart Budget :}
\begin{verbatim}
    +================================================================+
    |                 MY SMART BUDGET - ARCHITECTURE 3 TIERS        |
    +================================================================+

    +--------------------+    HTTPS/JWT      +--------------------+
    |   TIER 1           | <-------------->  |   TIER 2           |
    |   REACT FRONTEND   |   REST API        |   NODE.JS BACKEND  |
    |                    |   WebSocket       |                    |
    | +--------------+   |                   | +--------------+   |
    | | Dashboard    |   |                   | | Controllers  |   |
    | | Transactions |   |                   | | Services     |   |
    | | Budgets      |   |                   | | Middlewares  |   |
    | | Reports      |   |                   | | Validators   |   |
    | +--------------+   |                   | +--------------+   |
    +--------------------+                   +--------------------+
            |                                          |
            | Service Worker                           | SQL/NoSQL
            | (Offline sync)                           | Transactions
            v                                          v
    +--------------------+              +--------------------+
    |   CACHE LOCAL      |              |   TIER 3           |
    |   IndexedDB        |              |   DATA LAYER       |
    |   localStorage     |              |                    |
    +--------------------+              | +--------------+   |
                                       | | PostgreSQL   |   |
                                       | | MongoDB      |   |
                                       | | Redis Cache  |   |
                                       | +--------------+   |
                                       +--------------------+
                                       
    Flux de données typique:
    1. User -> React -> Validation locale
    2. React -> API REST -> JWT validation
    3. Controller -> Service -> Business logic
    4. Service -> Repository -> Database
    5. Database -> Repository -> Service
    6. Service -> Controller -> Response
    7. API -> React -> Update UI
    8. React -> Cache local (offline support)
\end{verbatim}
\end{exemple}

\subsection{Avantages de l'architecture pour My Smart Budget}

\textbf{Bénéfices spécifiques à l'application financière :}

\begin{exemple}
\textbf{Avantages de l'architecture :}
\begin{itemize}
    \item \textbf{Sécurité renforcée :} Isolation des données financières sensibles
    \item \textbf{Performance :} Cache Redis pour calculs répétitifs de budgets
    \item \textbf{Scalabilité :} Possibilité d'ajouter des workers pour les calculs
    \item \textbf{Offline first :} Synchronisation des transactions hors ligne
    \item \textbf{Audit trail :} MongoDB pour tracer toutes les opérations
    \item \textbf{Real-time :} WebSocket pour alertes de dépassement budget
    \item \textbf{Multi-device :} API REST permettant mobile et web
    \item \textbf{Conformité :} Architecture facilitant RGPD/PCI DSS
\end{itemize}
\end{exemple}

\begin{conseil}
\begin{itemize}
    \item Chiffrer toutes les données financières sensibles
    \item Implémenter une authentification à deux facteurs
    \item Logger toutes les opérations financières pour l'audit
    \item Utiliser des decimal types pour éviter les erreurs d'arrondi
    \item Mettre en cache les calculs de budget coûteux
    \item Prévoir la synchronisation offline pour mobile
    \item Implémenter des webhooks pour les banques
\end{itemize}
\end{conseil}

\begin{jury}
\begin{itemize}
    \item Comment sécurisez-vous les données financières des utilisateurs ?
    \item Quelle est votre stratégie de synchronisation offline ?
    \item Comment gérez-vous la précision des calculs monétaires ?
    \item Votre architecture supporte-t-elle plusieurs devises ?
    \item Comment tracez-vous les opérations pour l'audit ?
    \item Quelle est la latence de vos alertes de dépassement ?
    \item Comment gérez-vous les transactions récurrentes ?
\end{itemize}
\end{jury}

\section{Développement Frontend - Interface My Smart Budget}

Le frontend de My Smart Budget privilégie l'expérience utilisateur avec des tableaux de bord intuitifs, des visualisations de données claires et une navigation fluide entre les modules financiers.

\textbf{Technologies choisies :} React avec TypeScript pour la sécurité des types dans les calculs financiers, Material-UI pour une interface moderne et Chart.js pour les visualisations.

\textbf{Architecture des composants :} Organisation modulaire par domaine fonctionnel (transactions, budgets, rapports) avec des composants réutilisables pour l'affichage des montants.

\textbf{Accessibilité et UX :} Interface accessible WCAG AA avec support du mode sombre, navigation clavier complète et lecteurs d'écran pour l'inclusion financière.

\begin{exemple}
\textbf{Composant TransactionForm accessible :}
\begin{lstlisting}[language=JavaScript]
const TransactionForm = ({ onSubmit, categories, budget }) => {
  const [formData, setFormData] = useState({
    amount: '',
    category: '',
    description: '',
    type: 'expense',
    date: new Date()
  });
  
  const [errors, setErrors] = useState({});
  
  const validateAmount = (value) => {
    const amount = parseFloat(value);
    if (isNaN(amount) || amount <= 0) {
      return 'Le montant doit être supérieur à 0';
    }
    if (amount > 1000000) {
      return 'Le montant semble trop élevé';
    }
    return null;
  };
  
  const handleSubmit = async (e) => {
    e.preventDefault();
    const amountError = validateAmount(formData.amount);
    
    if (amountError) {
      setErrors({ amount: amountError });
      return;
    }
    
    // Vérification du budget disponible
    if (budget && parseFloat(formData.amount) > budget.remaining) {
      const confirm = await showConfirmDialog(
        'Dépassement de budget',
        `Cette transaction dépassera votre budget de ${
          parseFloat(formData.amount) - budget.remaining
        }€. Continuer ?`
      );
      
      if (!confirm) return;
    }
    
    await onSubmit(formData);
  };
  
  return (
    <form onSubmit={handleSubmit} aria-label="Ajouter une transaction">
      <FormControl fullWidth error={!!errors.amount}>
        <InputLabel htmlFor="amount-input">
          Montant *
        </InputLabel>
        <Input
          id="amount-input"
          type="number"
          step="0.01"
          value={formData.amount}
          onChange={(e) => setFormData({
            ...formData,
            amount: e.target.value
          })}
          startAdornment={<InputAdornment position="start">€</InputAdornment>}
          aria-describedby="amount-error"
          aria-required="true"
          inputProps={{
            'aria-label': 'Montant en euros',
            min: 0,
            max: 1000000
          }}
        />
        {errors.amount && (
          <FormHelperText id="amount-error" role="alert">
            {errors.amount}
          </FormHelperText>
        )}
      </FormControl>
      
      <FormControl fullWidth>
        <InputLabel id="category-label">Catégorie *</InputLabel>
        <Select
          labelId="category-label"
          value={formData.category}
          onChange={(e) => setFormData({
            ...formData,
            category: e.target.value
          })}
          aria-required="true"
        >
          {categories.map(cat => (
            <MenuItem key={cat.id} value={cat.id}>
              <ListItemIcon>{cat.icon}</ListItemIcon>
              <ListItemText>
                {cat.name}
                {cat.budget && (
                  <Typography variant="caption" color="textSecondary">
                    {' '}(Budget: {cat.budget.remaining}€ restants)
                  </Typography>
                )}
              </ListItemText>
            </MenuItem>
          ))}
        </Select>
      </FormControl>
      
      <Button
        type="submit"
        variant="contained"
        color="primary"
        fullWidth
        aria-label="Enregistrer la transaction"
      >
        Enregistrer
      </Button>
    </form>
  );
};
\end{lstlisting}
\end{exemple}

\begin{exemple}
\textbf{Dashboard avec graphiques :}
\begin{lstlisting}[language=JavaScript]
const BudgetDashboard = () => {
  const { budgets, transactions, loading } = useBudgetData();
  
  return (
    <Grid container spacing={3}>
      {/* Cartes de résumé */}
      <Grid item xs={12} md={4}>
        <Card>
          <CardContent>
            <Typography color="textSecondary" gutterBottom>
              Solde du mois
            </Typography>
            <Typography variant="h4" component="div">
              {formatCurrency(calculateBalance(transactions))}
            </Typography>
            <LinearProgress
              variant="determinate"
              value={calculateSavingsRate(transactions)}
              aria-label="Taux d'épargne"
            />
          </CardContent>
        </Card>
      </Grid>
      
      {/* Graphique des dépenses */}
      <Grid item xs={12} md={8}>
        <Card>
          <CardContent>
            <Typography variant="h6">
              Évolution des dépenses
            </Typography>
            <ResponsiveContainer width="100%" height={300}>
              <LineChart data={transactions}>
                <CartesianGrid strokeDasharray="3 3" />
                <XAxis dataKey="date" />
                <YAxis />
                <Tooltip formatter={formatCurrency} />
                <Legend />
                <Line
                  type="monotone"
                  dataKey="expenses"
                  stroke="#f44336"
                  name="Dépenses"
                />
                <Line
                  type="monotone"
                  dataKey="income"
                  stroke="#4caf50"
                  name="Revenus"
                />
              </LineChart>
            </ResponsiveContainer>
          </CardContent>
        </Card>
      </Grid>
    </Grid>
  );
};
\end{lstlisting}
\end{exemple}

\begin{conseil}
\begin{itemize}
    \item Formater les montants selon la locale de l'utilisateur
    \item Utiliser des couleurs contrastées pour l'accessibilité
    \item Implémenter le mode offline avec Service Workers
    \item Optimiser les rerenders avec React.memo
    \item Virtualiser les longues listes de transactions
    \item Précharger les données fréquemment consultées
\end{itemize}
\end{conseil}

\begin{jury}
\begin{itemize}
    \item Comment gérez-vous l'affichage de grandes listes de transactions ?
    \item Votre application fonctionne-t-elle hors ligne ?
    \item Quels sont vos scores Lighthouse pour une page de dashboard ?
    \item Comment optimisez-vous le chargement initial ?
    \item L'interface est-elle utilisable au clavier uniquement ?
    \item Comment gérez-vous les erreurs réseau ?
\end{itemize}
\end{jury}

\section{Développement Backend - API My Smart Budget}

L'API REST de My Smart Budget gère toute la logique financière avec une attention particulière à la sécurité, la validation des données et la performance des calculs budgétaires.

\begin{exemple}
\textbf{Structure backend My Smart Budget :}
\begin{verbatim}
src/
+-- controllers/
|   +-- authController.js          # Authentification/2FA
|   +-- transactionController.js   # CRUD transactions
|   +-- budgetController.js        # Gestion budgets
|   +-- analyticsController.js     # Rapports et stats
|   +-- bankController.js          # Sync bancaire
+-- services/
|   +-- transactionService.js      # Logique transactions
|   +-- budgetService.js           # Calculs budgets
|   +-- alertService.js            # Alertes dépassement
|   +-- bankSyncService.js         # Import bancaire
|   +-- mlService.js               # Catégorisation auto
+-- repositories/
|   +-- transactionRepository.js   # Accès DB transactions
|   +-- budgetRepository.js        # Accès DB budgets
|   +-- userRepository.js          # Accès DB users
+-- middleware/
|   +-- auth.js                    # JWT validation
|   +-- rateLimiter.js            # Protection API
|   +-- validator.js              # Validation données
|   +-- auditLogger.js            # Log des opérations
+-- utils/
|   +-- currencyConverter.js      # Conversion devises
|   +-- calculator.js             # Calculs financiers
|   +-- encryption.js             # Chiffrement données
\end{verbatim}

\textbf{Validation et sécurité :}
\begin{lstlisting}[language=JavaScript]
// Middleware de validation des transactions
const validateTransaction = [
  body('amount')
    .isFloat({ min: 0.01, max: 1000000 })
    .withMessage('Montant invalide')
    .customSanitizer(value => Math.round(value * 100) / 100),
  
  body('category')
    .isUUID()
    .withMessage('Catégorie invalide'),
  
  body('date')
    .isISO8601()
    .toDate()
    .custom(date => date <= new Date())
    .withMessage('La date ne peut pas être dans le futur'),
  
  body('description')
    .trim()
    .isLength({ max: 500 })
    .escape(), // Protection XSS
  
  (req, res, next) => {
    const errors = validationResult(req);
    if (!errors.isEmpty()) {
      return res.status(400).json({ 
        errors: errors.array(),
        message: 'Données invalides'
      });
    }
    next();
  }
];

// Protection contre les attaques par force brute
const loginLimiter = rateLimit({
  windowMs: 15 * 60 * 1000, // 15 minutes
  max: 5, // 5 tentatives max
  message: 'Trop de tentatives, réessayez dans 15 minutes',
  standardHeaders: true,
  legacyHeaders: false,
});
\end{lstlisting}
\end{exemple}

\section{Gestion des données - My Smart Budget}

La couche données de My Smart Budget utilise PostgreSQL pour les données transactionnelles critiques avec des contraintes strictes et MongoDB pour l'analytique avec des agrégations complexes.

\begin{exemple}
\textbf{Transactions avec Prisma :}
\begin{lstlisting}[language=JavaScript]
class BudgetTransactionService {
  async transferBetweenBudgets(userId, fromBudgetId, toBudgetId, amount) {
    // Transaction ACID pour garantir la cohérence
    return await prisma.$transaction(async (tx) => {
      // Vérifier le budget source
      const fromBudget = await tx.budget.findUnique({
        where: { id: fromBudgetId, userId }
      });
      
      if (!fromBudget || fromBudget.remaining < amount) {
        throw new Error('Fonds insuffisants');
      }
      
      // Débiter le budget source
      await tx.budget.update({
        where: { id: fromBudgetId },
        data: { 
          remaining: { decrement: amount },
          spent: { increment: amount }
        }
      });
      
      // Créditer le budget destination
      await tx.budget.update({
        where: { id: toBudgetId },
        data: { 
          remaining: { increment: amount },
          spent: { decrement: amount }
        }
      });
      
      // Logger la transaction
      await tx.budgetTransfer.create({
        data: {
          fromBudgetId,
          toBudgetId,
          amount,
          userId,
          timestamp: new Date()
        }
      });
      
      return { success: true, amount };
    });
  }
}
\end{lstlisting}

\textbf{Agrégations MongoDB pour analytics :}
\begin{lstlisting}[language=JavaScript]
// Analyse des habitudes de dépenses
async function analyzeSpendingPatterns(userId, period) {
  const pipeline = [
    // Filtrer par utilisateur et période
    {
      $match: {
        userId: new ObjectId(userId),
        date: {
          $gte: period.start,
          $lte: period.end
        },
        type: 'expense'
      }
    },
    // Grouper par jour de la semaine
    {
      $group: {
        _id: { $dayOfWeek: '$date' },
        avgAmount: { $avg: '$amount' },
        totalAmount: { $sum: '$amount' },
        count: { $sum: 1 },
        categories: { $addToSet: '$category' }
      }
    },
    // Ajouter des insights
    {
      $project: {
        dayOfWeek: '$_id',
        avgAmount: { $round: ['$avgAmount', 2] },
        totalAmount: { $round: ['$totalAmount', 2] },
        transactionCount: '$count',
        uniqueCategories: { $size: '$categories' },
        spendingIntensity: {
          $switch: {
            branches: [
              { case: { $lt: ['$avgAmount', 20] }, then: 'low' },
              { case: { $lt: ['$avgAmount', 50] }, then: 'medium' },
              { case: { $gte: ['$avgAmount', 50] }, then: 'high' }
            ]
          }
        }
      }
    },
    { $sort: { dayOfWeek: 1 } }
  ];
  
  return await db.collection('transactions').aggregate(pipeline).toArray();
}
\end{lstlisting}
\end{exemple}

\section{Sécurité et Conformité - My Smart Budget}

My Smart Budget implémente des mesures de sécurité strictes pour protéger les données financières sensibles des utilisateurs.

\begin{exemple}
\textbf{Chiffrement des données sensibles :}
\begin{lstlisting}[language=JavaScript]
// Chiffrement AES-256 pour les données bancaires
const crypto = require('crypto');

class EncryptionService {
  constructor() {
    this.algorithm = 'aes-256-gcm';
    this.key = Buffer.from(process.env.ENCRYPTION_KEY, 'hex');
  }
  
  encryptBankCredentials(credentials) {
    const iv = crypto.randomBytes(16);
    const cipher = crypto.createCipheriv(this.algorithm, this.key, iv);
    
    let encrypted = cipher.update(JSON.stringify(credentials), 'utf8', 'hex');
    encrypted += cipher.final('hex');
    
    const authTag = cipher.getAuthTag();
    
    return {
      encrypted,
      iv: iv.toString('hex'),
      authTag: authTag.toString('hex')
    };
  }
  
  decryptBankCredentials(encryptedData) {
    const decipher = crypto.createDecipheriv(
      this.algorithm,
      this.key,
      Buffer.from(encryptedData.iv, 'hex')
    );
    
    decipher.setAuthTag(Buffer.from(encryptedData.authTag, 'hex'));
    
    let decrypted = decipher.update(encryptedData.encrypted, 'hex', 'utf8');
    decrypted += decipher.final('utf8');
    
    return JSON.parse(decrypted);
  }
}
\end{lstlisting}

\textbf{Audit et conformité RGPD :}
\begin{lstlisting}[language=JavaScript]
// Service d'audit pour traçabilité
class AuditService {
  async logFinancialOperation(userId, operation, details) {
    await this.auditDb.collection('audit_logs').insertOne({
      userId,
      operation,
      details: this.sanitizeDetails(details),
      ip: this.getClientIp(),
      userAgent: this.getUserAgent(),
      timestamp: new Date(),
      sessionId: this.getSessionId(),
      risk_score: this.calculateRiskScore(operation, details)
    });
    
    // Détection d'anomalies
    if (this.isAnomalous(operation, details)) {
      await this.alertService.notifySecurityTeam({
        userId,
        operation,
        reason: 'Activité suspecte détectée'
      });
    }
  }
  
  async exportUserData(userId) {
    // Conformité RGPD - Export des données
    const userData = {
      profile: await this.userRepo.findById(userId),
      transactions: await this.transactionRepo.findByUser(userId),
      budgets: await this.budgetRepo.findByUser(userId),
      auditLogs: await this.getAuditLogs(userId)
    };
    
    return this.encryptionService.encryptForExport(userData);
  }
}
\end{lstlisting}
\end{exemple}

\begin{conseil}
\begin{itemize}
    \item Chiffrer toutes les données bancaires au repos
    \item Implémenter l'authentification à deux facteurs (2FA)
    \item Utiliser des tokens JWT avec expiration courte
    \item Logger toutes les opérations financières
    \item Scanner régulièrement les dépendances (npm audit)
    \item Implémenter le rate limiting sur toutes les routes
    \item Anonymiser les données dans les logs
    \item Prévoir l'export RGPD des données utilisateur
\end{itemize}
\end{conseil}

\begin{jury}
\begin{itemize}
    \item Comment protégez-vous les données bancaires ?
    \item Votre application est-elle conforme au RGPD ?
    \item Comment détectez-vous les transactions frauduleuses ?
    \item Quelle est votre stratégie de backup ?
    \item Comment gérez-vous les sessions utilisateur ?
    \item Avez-vous un plan de réponse aux incidents ?
    \item Comment auditez-vous les accès aux données sensibles ?
\end{itemize}
\end{jury}

\section{Liens utiles}

\begin{itemize}
    \item PCI DSS Compliance: \url{https://www.pcisecuritystandards.org/}
    \item RGPD pour applications financières: \url{https://www.cnil.fr/}
    \item Open Banking API: \url{https://www.openbanking.org.uk/}
    \item Plaid API (Agrégation bancaire): \url{https://plaid.com/docs/}
    \item Chart.js Documentation: \url{https://www.chartjs.org/docs/}
    \item JWT Best Practices: \url{https://tools.ietf.org/html/rfc8725}
    \item OWASP Top 10: \url{https://owasp.org/www-project-top-ten/}
\end{itemize}